\section{Scope}

This document sets out the PSLF model of Fervo Cape Station Phase~2 as delivered in
the two source files listed in Section~\ref{sec:refs}: a power-flow network case
(\texttt{.epc}, PSLF version~24.3.0)~[1] and a dynamic-data file (\texttt{.dyd})~[2].
The representation is split into a \emph{physical part} --- buses, lines, transformers,
generators and loads that define the steady-state network --- and a \emph{control
part} --- the machine, excitation, governor, stabilizer and protection models that
define the dynamic behaviour of each unit.

The plant comprises nine steam-turbine generators at 13.8\,kV, arranged in pairs on
four generator buses, each pair stepped up through its own 150\,MVA generator
step-up transformer (GSU) to a 345\,kV collector bus (GEN-TIE). The collector bus
connects over a 3.0\,mile line to the point of interconnection (POI), which is
represented as an infinite bus in the dynamic data. A 3.2\,mile line from GEN-TIE to
the Phase~1 main power transformer bus (HV-MPT-P1) is included in the network case
as the tie to Phase~1.

Two families of units appear in the data. Units~4, 5 and~6 are Turboden machines
rated 76.471\,MV$\cdot$A, and units~7 through~11 are Baker Hughes machines rated
71.13\,MV$\cdot$A. The two families share the same excitation, stabilizer and
protection settings but have different machine constants and turbine capabilities.
Every parameter value in the tables of this document is transcribed from the source
files by script, so the tables can be regenerated whenever the files change.
